% ============================================================
%  Chapitre 8 — Comparaison des formalismes
% ============================================================
\chapter{Comparaison des formalismes}

\section{Tableau comparatif}

\begin{table}[H]
\centering
\caption{Comparaison des cinq formalismes implémentés}
\label{tab:comparaison}
\renewcommand{\arraystretch}{1.4}
\begin{tabularx}{\linewidth}{lXXXXX}
\toprule
& \textbf{Modale} & \textbf{Défauts} & \textbf{Classique} & \textbf{Description} & \textbf{Rés. Sém.} \\
\midrule
\textbf{Monotone}
  & \textcolor{codegreen}{Oui}
  & \textcolor{red}{Non}
  & \textcolor{codegreen}{Oui}
  & \textcolor{codegreen}{Oui}
  & \textcolor{red}{Non*} \\

\textbf{Monde ouvert (OWA)}
  & Non
  & Non
  & Non
  & \textcolor{codegreen}{Oui}
  & Non \\

\textbf{Exceptions}
  & Non
  & \textcolor{codegreen}{Oui}
  & Non
  & Partiel
  & \textcolor{codegreen}{Oui} \\

\textbf{Quantificateurs}
  & Non
  & \textcolor{codegreen}{Oui (FOL)}
  & \textcolor{codegreen}{Oui (FOL)}
  & \textcolor{codegreen}{Oui ($\exists$, $\forall$)}
  & Non \\

\textbf{Mondes poss.}
  & \textcolor{codegreen}{Oui}
  & Non
  & Non
  & Non
  & Non \\

\textbf{Complexité}
  & PSPACE
  & $\Pi_2^P$
  & NP (SAT)
  & ExpTime
  & P \\

\textbf{Outil}
  & Tweety ML
  & Tweety RDL
  & SAT4J / Tweety
  & OWL API + HermiT
  & Tweety \\

\textbf{Standard}
  & Logiques K/T/S4/S5
  & Reiter 1980
  & ---
  & OWL 2 (W3C)
  & ---\\
\bottomrule
\multicolumn{6}{l}{\small * avec héritage à exceptions}
\end{tabularx}
\end{table}

\section{Expressivité et pouvoir de représentation}

\paragraph{Logique classique} C'est le socle — elle permet de vérifier la
cohérence et de dériver des conséquences, mais elle est monotone et ne permet
pas de représenter l'incertitude ou les exceptions.

\paragraph{Logique modale} Ajoute la notion de \emph{point de vue} et de
\emph{mondes possibles} : on peut exprimer ce qui est vrai dans tous les scénarios
($\nec$) ou dans au moins un ($\pos$). Particulièrement adaptée au raisonnement
épistémique (ce qu'un agent \emph{sait} ou \emph{croit}).

\paragraph{Logique des défauts} Permet le raisonnement avec des \emph{règles
générales avec exceptions}. Le résultat (extensions) n'est pas unique en cas
de conflits — deux agents raisonnables peuvent tirer des conclusions différentes.
C'est le formalisme le plus adapté au raisonnement du sens commun.

\paragraph{Logique de description} Offre un cadre rigoureux et standardisé
(OWL~2) pour la modélisation de hiérarchies de concepts complexes.
Le raisonneur HermiT garantit la complétude et la correction.
L'OWA la distingue des autres formalismes.

\paragraph{Réseaux sémantiques} Simple et intuitif, particulièrement adapté
à la représentation de hiérarchies et à l'héritage de propriétés. Moins
expressif que la logique de description mais plus lisible pour des non-experts.

\section{Traitement de la non-monotonie}

\begin{table}[H]
\centering
\caption{Mécanismes de non-monotonie par formalisme}
\begin{tabularx}{\linewidth}{lXX}
\toprule
\textbf{Formalisme} & \textbf{Mécanisme} & \textbf{Illustration dans le projet} \\
\midrule
Logique des défauts
  & Ajout d'un fait dans $W$ peut invalider une conclusion d'extension
  & SigneLabel(flenn) invalide LibreDroits \\
Réseaux sémantiques
  & Redéfinition locale d'une propriété héritée
  & Exception à l'héritage \\
Logique de description
  & Non-monotonie simulée via deux ontologies distinctes
  & signe\_chez(flenn, universal) invalide l'assertion ArtisteLibreDroits \\
Logique modale
  & Pas de non-monotonie intrinsèque
  & --- \\
Logique classique
  & Pas de non-monotonie intrinsèque
  & --- \\
\bottomrule
\end{tabularx}
\end{table}

\section{Quand utiliser quel formalisme ?}

\begin{description}
  \item[Logique classique] Vérification formelle, contraintes dures,
        satisfiabilité, inférence déterministe. Idéal pour les systèmes
        où toutes les règles sont certaines et sans exception.

  \item[Logique modale] Raisonnement sur des états alternatifs du monde :
        croyances, obligations, permissions, temps. Adapté aux systèmes
        multi-agents, au droit, et à la vérification de protocoles.

  \item[Logique des défauts] Raisonnement du sens commun avec exceptions et
        révision de croyances. Adapté aux systèmes experts où les règles ont
        des exceptions et les informations peuvent être incomplètes.

  \item[Logique de description] Modélisation formelle de domaines structurés
        (ontologies). Adapté au Web sémantique, aux bases de connaissances
        médicales, juridiques ou scientifiques (interopérabilité via OWL).

  \item[Réseaux sémantiques] Représentation simple et visuelle de hiérarchies.
        Adapté aux applications nécessitant une sémantique lisible par des
        humains, comme les systèmes de recommandation ou les moteurs de recherche.
\end{description}

\section{Cohérence entre les domaines d'application}

Un point fort de ce projet est l'application systématique des cinq formalismes
aux \emph{mêmes domaines}. Cela permet d'observer comment chaque formalisme
capture différents aspects du même problème.

\begin{tcolorbox}[resultbox, title={Exemple : Flenn (industrie musicale)}]
\begin{itemize}
  \item \textbf{Défauts :} avant SigneLabel $\to$ extension \{LibreDroits\} ;
        après SigneLabel $\to$ deux extensions (conflit $d_1$/$d_2$).
  \item \textbf{DL :} avant signe\_chez $\to$ ArtisteLibreDroits=VRAI ;
        après signe\_chez $\to$ ArtisteLibreDroits=FAUX.
  \item \textbf{Modale :} le monde $w_3$ (artiste indépendant) est accessbile
        depuis $w_1$ (label) : $\pos\,\text{artiste\_independant}$ en $w_1$.
  \item \textbf{Classique :} clause $C_1 = \lnot\text{SigneLabel} \lor
        \lnot\text{LibreDroits}$ — cohérence vérifiée par SAT4J.
\end{itemize}
\end{tcolorbox}
